\documentclass[10pt,a4paper]{amsart}
\usepackage[margin=19mm]{geometry}
\usepackage{amsmath,amssymb,mathtools}
\usepackage[scr=rsfs,cal=euler]{mathalfa}
\usepackage{xfrac}
\usepackage[unicode,hidelinks,bookmarks=false]{hyperref}
\newcommand{\ie}{i.e.\ }
\newcommand{\cf}{cf.\ }
\newcommand{\resp}{respectively\ }
\newcommand{\Subsection}{\subsection}
\providecommand{\nobreakdash}{\nobreak\hbox{-}}
\setlength{\emergencystretch}{2em}
\multlinegap0pt
% ---------------------------------------------------------------------------
% Editorial AMS wrapper prelude. The theorem environments and the mathematical macros below are
% transcribed from the paper's own preamble (cached-originals/source0.tex lines 18--102); the AMS amsart
% class, the named format packages of the pinned TeX Live runtime and this editorial formatting are not
% part of the author's text. No mathematics is changed.
% ---------------------------------------------------------------------------
\usepackage{tikz-cd}
\newcommand{\bluebf}[1]{{\color{blue} \bf #1}}
\newcommand{\blue}[1]{{\color{blue} #1}}
\newcommand{\red}[1]{{\color{red} #1}}






\newtheorem{thm}{Theorem}
\newtheorem{prop}{Proposition}
\newtheorem{lem}{Lemma}
\newtheorem{cor}{Corollary}

\theoremstyle{definition}
\newtheorem{defn}{Definition}
\newtheorem{rem}{Remark}


\newtheorem{example}{Example}


\DeclarePairedDelimiter\abs{\lvert}{\rvert}
\DeclareMathOperator{\LT}{LT}
\DeclareMathOperator{\Lex}{Lex}
\DeclareMathOperator{\Image}{Im}
\DeclareMathOperator{\Ch}{Ch}
\DeclareMathOperator{\at}{at}
\DeclareMathOperator{\Con}{Con}
\DeclareMathOperator{\Q}{Q}
\DeclareMathOperator{\Fil}{Fil}
\DeclareMathOperator{\Ker}{Ker}
\DeclareMathOperator{\Formulas}{Form}
\DeclareMathOperator{\Pro}{Pro-}
\DeclareMathOperator{\Sub}{Sub}
\DeclareMathOperator{\Norm}{Norm}
\DeclareMathOperator{\Seq}{Seq}
\DeclareMathOperator{\internal}{\delta^i}
\DeclareMathOperator{\external}{\delta^e}
\DeclareMathOperator{\confl}{e}
\DeclareMathOperator{\height}{ht}
\DeclareMathOperator{\width}{wt}
\DeclareMathOperator{\confluent}{e}
\DeclareMathOperator{\leaf}{L}
\DeclareMathOperator{\Succ}{Succ}
\DeclareMathOperator{\Rank}{Rk}
\newcommand{\Ord}{\textup{Ord}}
\newcommand{\heart}{\ensuremath\heartsuit}
\newcommand{\KFr}{\textup{\textbf{KFr}}}
\newcommand{\PreO}{\textup{\textbf{PreO}}}
\newcommand{\Pos}{\textup{\textbf{Pos}}}
\newcommand{\GLLin}{\textup{\textbf{GL-Lin}}}
\newcommand{\MA}{\textup{\textbf{MA}}}
\newcommand{\BA}{\textup{\textbf{BA}}}
\newcommand{\CMA}{\textup{\textbf{CMA}}}
\newcommand{\CAMA}{\textup{\textbf{CAMA}}}
\newcommand{\CABA}{\textup{\textbf{CABA}}}
\newcommand{\fin}{\textit{fin}}
\newcommand{\lf}{\textit{lf}}
\newcommand{\id}{\textup{id}}
\newcommand{\op}{\textup{op}}
\newcommand{\Set}{\textup{\textbf{Set}}}
\newcommand{\Stone}{\textup{\textbf{Stone}}}
\newcommand{\concat}{\mathbin{\raisebox{0.9ex}{$\smallfrown$}}}
\newcommand{\Prop}{\textup{Prop}}
\newcommand{\LAlg}{L\textup{Alg}}


\begin{document}
\title{A Proof Theory for Profinite Modal Algebras}
\author{Matteo De Berardinis}
\address{University of Amsterdam}
\author{Silvio Ghilardi}
\address{Universit\`a degli Studi di Milano}
\date{Source: arXiv:2507.06007v4; distributed under CC BY 4.0}
\maketitle
\noindent\emph{Source, version and licence.} \emph{A Proof Theory for Profinite Modal Algebras}, by Matteo De Berardinis and Silvio Ghilardi. The source of this editable unit is the e-print of \textbf{version v4} of arXiv:2507.06007, https://arxiv.org/abs/2507.06007v4, whose material is distributed under the Creative Commons Attribution 4.0 International licence, http://creativecommons.org/licenses/by/4.0/, as recorded in the arXiv licence metadata for this paper and shown on that abstract page. The whole of the body below is version v4, the newest version of the paper. Full rights notice: licenses/arxiv-2507.06007-CC-BY-4.0.txt.\par\smallskip
\noindent\textit{Editorial scope.} Counted unit: original Theorem 12 of the article, the statement carried by the source environment \texttt{thm} with source label \texttt{thm:craig}, cached-originals/source0.tex lines 2278--2280, printed as \textbf{Theorem 12} exactly as in the article: Local Craig's Theorem holds for $\vdash_L$ if and only if the forgetful functor $L\KFr_\lf\longrightarrow\Set$ preserves weak pullbacks. It is delivered together with the author's own complete proof as the single primary proof body, source0.tex lines 2282--2323 (42 lines): the \texttt{proof} environment that immediately follows the statement, that reduces the superamalgamation property of a commutative square of profinite algebras to the corresponding condition on the dual square of locally finite Kripke frames, plays the inverse-image adjunctions to reduce it to a condition on singletons, and identifies the result with the weak-pullback condition in $\Set$. No proof marker is added and no mathematics is written, repaired or re-derived; the author's notation and the printed slips of the source are kept literal. Necessary complete context is retained uncounted, in source order: source0.tex lines 268--280, the definition of a locally finite Kripke frame together with the notation $L\KFr_\lf$ that the counted statement and proof use; lines 2070--2072, Theorem 10, the global Craig/amalgamation equivalence of which the counted theorem is the local counterpart; and lines 2272--2274, Proposition 11, the superamalgamation characterization that the counted proof invokes at its first line. The article's own numbering is reproduced by explicit counter settings: this article declares each theorem-like environment with its own global counter (\texttt{\textbackslash newtheorem\{thm\}\{Theorem\}} and likewise for prop, lem, cor, defn, rem and example), and every one of those counters is set before each retained passage, so the counted statement prints as Theorem 12, Theorem 10 as 10 and Proposition 11 as 11, exactly as in the article. No \texttt{\textbackslash cite} occurs in any retained passage: the closure of the retained passages over references and citations was computed, it contains no reference and no citation at all, so this unit delivers no bibliography entry and the article's own bibliography data file \texttt{refs.bib} is neither spliced into the bundled source nor redistributed. The retained passages contain no figure environment and no graphics-inclusion command, so no artwork is redistributed; the commutative diagrams of the counted proof are drawn with the named \texttt{tikz-cd} package of the pinned TeX Live runtime, exactly as the source does. The source's own class is AMS \texttt{amsart}; the e-print ships the authors' own style files \texttt{modalops.sty} and \texttt{prftree.sty}, neither of which is shipped, loaded or redistributed here (the mathematical macros the retained passages need are transcribed from the paper's own preamble instead), and no publisher class is shipped, used or redistributed. Every declared body of this unit is an exact, unmodified span of source0.tex: no body is a bounded passage comparison, \texttt{omit} was not used anywhere, and consequently no fidelity receipt is asserted in delivery-evidence.json. The complete concatenated original source is bundled as sources/181224/source0.tex. Omitted from this editable unit and therefore absent from it are source0.tex lines 1--267; 281--2069; 2073--2271; 2275--2277; 2281--2281; 2324--2528, that is the article's own preamble and title block, the introduction, the whole infinitary-calculus development, the r-regularity and exactness sections, the remaining interpolation results including the proof of Theorem 10, and the article's bibliography data, all of which remain present in the bundled source but are not delivered here. The AMS \texttt{amsart} wrapper, the named format packages of the pinned TeX Live runtime and this editorial source, licence and scope paragraph are editorial formatting only.\par\smallskip
\setcounter{section}{2}
\setcounter{equation}{0}
\setcounter{cor}{0}
\setcounter{defn}{0}
\setcounter{example}{0}
\setcounter{lem}{0}
\setcounter{prop}{0}
\setcounter{rem}{0}
\setcounter{thm}{1}
\begin{defn}
A Kripke frame $W$ is said to be \emph{locally-finite} if, for all $w \in W$,
$$w^* \coloneqq \{w' \in W\ \vert\ \exists w_0,\dots,w_n \in W \text{ such that } w=w_0 \prec \dots \prec w_n=w'\}$$
is a finite set. In other words, fixing any starting point $w \in W$ and taking finite walks along the relation $\prec$, one can visit only finitely many $w' \in W$.
\end{defn}
We denote by $L\KFr_\lf$ the full subcategory of $L\KFr$ consisting of its locally-finite Kripke frames.

%\begin{rem}
%The class $L\KFr_\lf$ ($L\KFr_\fin$) is closed under (finite) disjoint unions, p-morphic images and generated subframes.
%\end{rem}

%We give some results whose proofs can be found in \cite{MDB-SG}.


\setcounter{section}{5}
\setcounter{equation}{0}
\setcounter{cor}{2}
\setcounter{defn}{7}
\setcounter{example}{2}
\setcounter{lem}{2}
\setcounter{prop}{10}
\setcounter{rem}{7}
\setcounter{thm}{9}
\begin{thm}\label{thm:craig=amalg}
Craig's Theorem holds for $\vdash_L$ if and only if $\Pro L\MA_\fin$ has the amalgamation property, if and only if $L\KFr_\lf$ is r-regular.
\end{thm}

\setcounter{section}{5}
\setcounter{equation}{0}
\setcounter{cor}{2}
\setcounter{defn}{8}
\setcounter{example}{2}
\setcounter{lem}{3}
\setcounter{prop}{10}
\setcounter{rem}{8}
\setcounter{thm}{11}
\begin{prop}
Local version of Craig's Theorem holds for $\vdash_L$ if and only if $\Pro L\MA_\fin$ has the superamalgamation property.
\end{prop}

\setcounter{section}{5}
\setcounter{equation}{0}
\setcounter{cor}{2}
\setcounter{defn}{8}
\setcounter{example}{2}
\setcounter{lem}{3}
\setcounter{prop}{11}
\setcounter{rem}{8}
\setcounter{thm}{11}
\begin{thm}\label{thm:craig}
Local Craig's Theorem holds for $\vdash_L$ if and only if the forgetful functor $L\KFr_\lf \longrightarrow \Set$ preserves weak pullbacks.
\end{thm}

\setcounter{section}{5}
\setcounter{equation}{0}
\setcounter{cor}{2}
\setcounter{defn}{8}
\setcounter{example}{2}
\setcounter{lem}{3}
\setcounter{prop}{11}
\setcounter{rem}{8}
\setcounter{thm}{12}
\begin{proof}
First, let us observe that $L$ has (SAP) if and only if all pushouts in $\Pro L\MA_\fin$ have the superamalgamation property, if and only if all weak pushouts in $\Pro L\MA_\fin$ have the superamalgamation property (use the universal properties to establish these equivalences). Thus it is sufficient to show that a commutative square of profinite algebras 
    \[\begin{tikzcd}
	A & B \\
	C & D
	\arrow["k", from=1-1, to=1-2]
	\arrow["g"', from=1-1, to=2-1]
	\arrow["f", from=1-2, to=2-2]
	\arrow["h"', from=2-1, to=2-2]
    \end{tikzcd}\]
    has the superamalgamation property if and only if the dual square 
    \[\begin{tikzcd}
	A^* & B^* \\
	C^* & D^*
	\arrow["k^*", from=1-2, to=1-1]
	\arrow["g^*"', from=2-1, to=1-1]
	\arrow["f^*", from=2-2, to=1-2]
	\arrow["h^*"', from=2-2, to=2-1]
    \end{tikzcd}\]
     of locally finite Kripke frames becomes a weak pullback once the forgetful functor is applied to it.
     This is easily established as follows. The (SAP) property for the dualized square means that, for all $c\subseteq C^*, b\subseteq B^*$, we have
     $$
     (h^*)^{-1}(c) \subseteq (f^*)^{-1}(b)~~\Rightarrow~~
     \exists a\subseteq A^*~(c\subseteq  (g^*)^{-1}(a)~\&~
      (k^*)^{-1}(a)\subseteq b)~~.
     $$
     Playing with the fact that inverse image has both a left and a right adjoint, this is the same as asking that, for all $c\subseteq C^*, b\subseteq B^*$, we have 
     $$
     \exists_{f^*}(h^*)^{-1}(c)\subseteq b 
     ~~\Rightarrow~~(k^*)^{-1}\exists_{g^*}(c)\subseteq b 
     $$
     (here $\exists_{f^*}$ and $\exists_{g^*}$ are the left adjoints). Since the latter has to hold for all $b$, it reduces to requiring 
     $$
     (k^*)^{-1}\exists_{g^*}(c) \subseteq  \exists_{f^*}(h^*)^{-1}(c)
     $$
     for all $c\subseteq C^*$. In turn, such condition is verified if and only if it is verified when $c$ is a singleton, i.e.\ if and only if
     $$
     \forall x\in C^*\, \forall y\in B^*\,(k^*(y)=g^*(x) 
     ~\Rightarrow~\exists z\in D^* (h^*(z)=x ~\&~f^*(z)=y))
     $$
     The latter  is precisely the condition for a commutative square in $\Set$ to be a weak pullback.
\end{proof}

\end{document}
